% ============================================================
\section{End (Endomaps)}
As construções na categoria de endomapas focam-se na preservação da evolução de estados. O foco principal é garantir que as novas estruturas criadas mantenham uma dinâmica válida, onde a "comutatividade" da evolução do sistema original é respeitada.

\subsection{Fundamentação Teórica}
As principais construções universais implementadas para esta categoria são:
\begin{itemize}
    \item \textbf{Produtos:} O produto $X \times Y$ atua em paralelo com o endomapa produto definido por $\alpha \times \beta (x, y) = (\alpha(x), \beta(y))$.
    \item \textbf{Equalizadores:} O subconjunto de estados invariantes face aos morfismos concorrentes, munido da restrição da dinâmica original.
\end{itemize}

\subsection{Implementação Computacional}
Em vez do produto de Kronecker, iteramos e indexamos linearmente a combinação de vetores para construir o endomapa produto diretamente num vetor unidimensional.

\begin{lstlisting}[caption={Construção do Produto Paralelo de Endomapas}]
function P = endo_product(EA, EB)
    % EA e EB são vetores de endomapas
    nA = length(EA);
    nB = length(EB);
    P_vector = zeros(1, nA * nB);
    
    for i = 1:nA
        for j = 1:nB
            % Cálculo do índice no produto cartesiano
            idx = (i-1)*nB + j;
            % Onde (EA(i), EB(j)) cai no novo vetor referencial
            val = (EA(i)-1)*nB + EB(j);
            P_vector(idx) = val;
        end
    end
    
    P.vector = P_vector;
    P.n_elements = nA * nB;
end
\end{lstlisting}